\par\SourceLeaf{30}
{\rightskip=0pt plus 1em Der eigenthümliche Charakter des Princips der virtuellen Geschwin\-digkeiten besteht darin, dass es eine allgemeine Formel zur Auflösung aller statischen Aufgaben, und so der Stellvertreter aller andern Principe ist, ohne jedoch das Creditiv dazu so unmittelbar aufzuweisen, dass es sich, so wie es nur ausgesprochen wird, schon von selbst als plausibel empföhle.\par}

In \textit{dieser} Beziehung scheint das Princip, welches ich hier aufstellen werde, den Vorzug zu haben: es hat aber auch noch den zweiten, dass es das Gesetz der Bewegung und der Ruhe auf ganz gleiche Art in grösster Allgemeinheit umfasst. So sehr es in der Ordnung ist, dass bei der allmäligen Ausbildung der Wissenschaft und bei der Belehrung des Individuum das Leichtere dem Schwereren, das Einfachere dem Verwickeltern, das Besondere dem Allgemeinen vorangeht, so fordert doch der Geist, einmal auf dem höhern Standpunkte angelangt, den umgekehrten Gang, wobei die ganze Statik nur als ein ganz specieller Fall der Mechanik erscheine. Selbst der oben erwähnte Geometer scheint darauf Werth zu legen, indem er als einen Vorzug des Princips der kleinsten Wirkung ansieht, dass es das Gleichgewicht und die Bewegung zugleich umfasse, wenn man jenes so ausdrücke, dass die lebendigen Kräfte bei beiden Kleinste seien, eine Bemerkung, die doch mehr witzig als wahr zu sein scheint, da das Minimum in beiden Fällen in ganz verschiedener Beziehung Statt findet.

Das neue Princip ist nun folgendes:

\textit{Die Bewegung eines Systems materieller, auf was immer für eine Art unter sich verknüpfter Punkte, deren Bewegungen zugleich an was immer für äussere Beschränkungen gebunden sind, geschieht in jedem Augenblick in möglich grösster Übereinstimmung mit der freien Bewegung, oder unter möglich kleinstem Zwange, indem man als Maass des Zwanges, den das ganze System in jedem Zeittheilchen erleidet, die Summe der Producte aus dem Quadrate der Ablenkung jedes Punkts von seiner freien Bewegung in seine Masse betrachtet.}

Es seien $m,m',m''$ u. s. w. die Massen der Punkte; $a,a',a''$ u. s. w. ihre Plätze zur Zeit $t$; $b,b',b''$ u. s. w. die Plätze, welche sie, nach dem unendlich
